\documentclass[11pt]{article}
\usepackage[a4paper,margin=27mm,headheight=14pt]{geometry}
\usepackage[T1]{fontenc}
\usepackage{lmodern,microtype,amsmath,amssymb,amsthm,mathtools,booktabs}
\usepackage[colorlinks=true,linkcolor=blue,citecolor=blue,urlcolor=blue]{hyperref}
\newtheorem{theorem}{Theorem}[section]
\newtheorem{lemma}[theorem]{Lemma}
\newtheorem{proposition}[theorem]{Proposition}
\newtheorem{corollary}[theorem]{Corollary}
\theoremstyle{remark}\newtheorem{remark}[theorem]{Remark}
\hypersetup{pdftitle={A second sharp jump in sparse p-adic root counts},pdfauthor={ProveIt research note}}
\newcommand{\Z}{\mathbb Z}\newcommand{\Fp}{\mathbb F_p}\newcommand{\Zp}{\mathbb Z_p}\newcommand{\Qp}{\mathbb Q_p}\newcommand{\Cp}{\mathbb C_p}
\newcommand{\vp}{v_p}\newcommand{\Ncirc}{N_p^\circ}\newcommand{\dd}{\mathrm d}
\setlength{\emergencystretch}{2em}
\title{A second sharp jump in sparse $p$-adic root counts\\[3pt]\large An exact threshold at $2p+4$ monomials}
\author{Research note prepared for the ProveIt project}
\date{27 November 2026}
\begin{document}
\maketitle
\begin{abstract}
For a prime $p$, let $N_p^\circ(t)$ be the greatest number of roots, with multiplicities, that a nonzero polynomial with at most $t$ monomials can have in $v_p(X-1)>1/(p-1)$, with coefficients in a finite extension of $\mathbb Q_p$. For every prime $p\ge5$ we prove
\[
N_p^\circ(2p+2)=p^2+\frac{p-1}{2},\qquad
N_p^\circ(2p+3)=p^2+\frac{p+1}{2},\qquad
N_p^\circ(2p+4)=p^2+p.
\]
Thus one additional monomial at the last threshold increases the sharp bound by $(p-1)/2$. The last value is attained by an integer polynomial of degree $p^2+p$, with exactly $2p+4$ nonzero coefficients, all of whose roots are distinct and belong to the stated disk. The support uses three residue classes with multiplicities $(p+2,p+1,1)$. The proof reduces high moments to a finite algebra over $\mathbb F_p$ and supplies explicit quadratic and cubic reduction identities. It extends the latest ProveIt sparse-root report, which proves the range through $2p+1$. All theorems here are unconditional.
\end{abstract}

\section{The problem and the new threshold}
Normalize $v_p(p)=1$ and work in $\overline{\mathbb Q}_p\subset\mathbb C_p$. Put
\[
\mathcal D^\circ=\{X:v_p(X-1)>1/(p-1)\},\qquad
\mathcal D(r)=\{X:v_p(X-1)\ge r\}\quad(r>1/(p-1)).
\]
The number $N_p^\circ(t)$ is the supremum of the numbers of roots in $\mathcal D^\circ$ of nonzero polynomials with at most $t$ monomials. Multiplicities are counted, exponents are distinct nonnegative integers, and coefficients may belong to any finite extension of $\mathbb Q_p$. The maxima below are also attained when roots must be distinct and coefficients must be integers. These are geometric roots in the algebraic closure; they are not asserted to lie in $\mathbb Q_p$ itself.

The merged ProveIt article asks for sharp bounds beyond $p+2$ terms in question \texttt{smd:q:roots} \cite{article}. The recent report \cite{transport} proves exact transport of root depths using finitely many exponent digits, a finite sharp algorithm for arbitrary support size, and
\begin{equation}\label{eq:old}
N_p^\circ(t)=p^2+\left\lfloor\frac{t-p}{2}\right\rfloor-1
\qquad(p+2\le t\le2p+1).
\end{equation}
We determine the next three values. The main new feature is the third residue class: a degree-two moment involving two distinct positive base-$p$ digits becomes independent of all earlier moments of at most that digit sum.

\begin{theorem}[The second jump]\label{thm:main}
For every prime $p\ge5$,
\begin{align}
 N_p^\circ(2p+2)&=p^2+(p-1)/2,\label{eq:first}\\
 N_p^\circ(2p+3)&=p^2+(p+1)/2,\label{eq:second}\\
 N_p^\circ(2p+4)&=p^2+p.\label{eq:third}
\end{align}
Each maximum is attained by an integer polynomial with simple roots. For \eqref{eq:third}, one may require exactly the support
\begin{equation}\label{eq:support}
 E=\{pj:0\le j\le p+1\}\ \cup\ \{1+pj:0\le j\le p\}\ \cup\ \{2\}.
\end{equation}
It has $2p+4$ elements and largest exponent $K=p^2+p$. Every root of this polynomial lies in $\mathcal D^\circ$, and all $K$ roots can already be required to lie in $\mathcal D(r)$ whenever
\begin{equation}\label{eq:depth}
 r=\frac1{p-1}+\delta,\qquad 0<\delta K<\frac1{p-1},\qquad \delta\in\mathbb Q.
\end{equation}
\end{theorem}

For example, the successive maxima at $t=12,13,14$ for $p=5$ are $27,28,30$; at $t=16,17,18$ for $p=7$ they are $52,53,56$. Formula \eqref{eq:old} therefore extends for two further terms, then fails at the next term. The new jump has size $(p-1)/2$, which is unbounded as $p$ grows.

\section{From zeros to reduced moments}
We give the moment reduction explicitly, so that the proof does not depend on the arbitrary-order transport theorem in \cite{transport}.

Let $P(X)=\sum_{e\in E}c_eX^e$, $|E|=t$, and write
\[
 P(\exp Z)=\sum_{j\ge0}\frac{M_j}{j!}Z^j,
 \qquad M_j=\sum_{e\in E}c_e e^j.
\]
The exponential and logarithm are inverse analytic maps between $v_p(Z)>1/(p-1)$ and $\mathcal D^\circ$, and preserve the valuation of $Z$ and $X-1$. Consequently they preserve zeros and their multiplicities on all the disks under discussion. Let
\[
 g(T)=\prod_{e\in E}(T-e),\qquad \Lambda=\Zp[T]/g(T).
\]
The $t$ initial moments specify a linear functional $L$ on the power basis $1,T,\ldots,T^{t-1}$ of $\Lambda$, and $M_j=L(T^j)$ for every $j$. Division by the monic polynomial $g$ is integral. After scaling the nonzero polynomial, we may therefore assume that all $M_j$ are integral in the coefficient field and at least one of $M_0,\ldots,M_{t-1}$ is a unit.

Write $S_p(j)$ for the base-$p$ digit sum. For $r=1/(p-1)+\delta$, Legendre's factorial formula gives
\begin{equation}\label{eq:weight}
 v_p(M_j/j!)+rj=v_p(M_j)+\frac{S_p(j)}{p-1}+\delta j.
\end{equation}
On every rational closed disk this quantity tends to infinity. The Newton polygon or Weierstrass preparation theorem says that the number of zeros, counted with multiplicity, is the last index minimizing \eqref{eq:weight}.

\begin{lemma}[Eligible high indices]\label{lem:eligible}
Suppose $p\ge5$ and $2p+2\le t\le2p+4$. A minimizing index $k\ge t$ satisfies
\begin{equation}\label{eq:eligible}
 S_p(k)\le p-1,\qquad v_p(M_k)<1.
\end{equation}
Let $\nu_c=\#\{e\in E:e\equiv c\pmod p\}$ and
\[
 R=\Fp[T]\big/\prod_{c\in\Fp}(T-c)^{\nu_c}.
\]
If $T^k$ in $R$ is a linear combination of powers $T^j$ with $j<k$ and $S_p(j)\le S_p(k)$, then $k$ cannot be a minimizing index.
\end{lemma}
\begin{proof}
Here $t\le3p-1$, so the largest digit sum among $0\le j<t$ is $p$. A unit initial moment has weight at most $p/(p-1)+\delta(t-1)$. For $k\ge t$, either $S_p(k)\ge p$ or $v_p(M_k)\ge1$ makes its weight strictly larger (in the latter case $S_p(k)\ge1$). This proves \eqref{eq:eligible}.

Lift a proposed relation to the integral power basis of $\Lambda$. Applying the integral functional $L$ gives the corresponding relation among moments modulo $p$ times the valuation ring. Since $v_p(M_k)<1$, at least one moment on the right has valuation at most $v_p(M_k)$. Its smaller index and no larger digit sum make its weight strictly smaller in \eqref{eq:weight}. A zero class is the same contradiction with an empty right side.
\end{proof}

\section{The finite-algebra reduction}
Set
\[
 x=T,\qquad y=T^p,\qquad z=T^{p^2}\quad\text{in }R.
\]
Since every $\nu_c\le2p+4<p^2$, on the local factor at $c$, with coordinate $u=T-c$, these elements are
\begin{equation}\label{eq:local}
 x=c+u,\qquad y=c+u^p,\qquad z=c.
\end{equation}
In particular $T^{p^a}=z$ for every $a\ge2$. A power with digit sum $s\le p-1$ can first be written as a monomial $x^a y^b z^c$ of total degree $s$ and no larger numerical exponent. If a base-$p$ digit above the $p^2$ place was present, this already strictly decreases the exponent.

Assign the monomial $x^a y^b z^c$ the numerical index $a+pb+p^2c$. All reductions below replace a monomial by a linear combination of strictly smaller indices and do not increase total degree. Multiplying a relation by another monomial preserves these properties. Since all relevant total degrees are at most $p-1$, no carrying between the three exponents is needed; more generally carrying could only decrease digit sum.

Call a residue \emph{thick} if $\nu_c>p$. There are at most two thick residues. The proof uses only a short list of possible multiplicity patterns.

\begin{lemma}[The cubic cutoff]\label{lem:cubic}
If $\sum_c\nu_c=2p+4$, every $T^k$ with $S_p(k)\le p-1$ and $k>p^2+p$ belongs to the span of $T^j$ with $j<k$ and $S_p(j)\le S_p(k)$.
\end{lemma}
\begin{proof}
We show that the leading monomials $z^2$, $xyz$, and $y^2z$ can be eliminated by relations of total degree at most two, three, and three, respectively. In several cases the stronger elimination of $yz$ is available. The remaining monomials are $x^ay^b$, $zx^a$, and possibly $yz$; of degree at most $p-1$, their indices are at most $p^2+p$.

\emph{No thick residue, or just one occupied residue.} With no thick residue, $y=z$ and the variable $z$ is eliminated. If just one residue $c$ is occupied, $z=c$ and the same conclusion holds.

\emph{One thick residue of multiplicity at most $2p$.} Let its residue be $c$ and put $v=y-z$. Then $v^2=0$ on every local factor, and $(z-c)v=0$, since $v$ is supported only at the thick residue. Their expansions give
\begin{equation}\label{eq:one-moderate}
 yz=y^2-cy+cz,\qquad z^2=y^2-2cy+2cz.
\end{equation}
These eliminate $yz$ and $z^2$.

\emph{One thick residue of multiplicity greater than $2p$.} The other multiplicities have sum $m\le3$. Translate $x,y,z$ by the thick residue, so that it becomes zero. Translating back later adds only proper divisors of the displayed leading monomials, with lower degree and index. On the thick factor $z=0$; on every other factor $y=z$, since each of its multiplicities is at most three and $p\ge5$. Thus $z^2=yz$.

Let $R_{\rm other}$ be the product of the other local factors. It is generated by $x$ and has dimension $m\le3$. Choose polynomials $A(x),B(x)$ of degrees less than $m$ representing respectively $xy$ and $y^2$ in this algebra, by division by its monic defining polynomial. Because multiplication by $z$ kills the thick factor,
\begin{equation}\label{eq:one-large}
 xyz=zA(x),\qquad y^2z=zB(x).
\end{equation}
These are the required cubic reductions; their right sides have degree at most three and index at most $p^2+2$.

\emph{Two thick residues and only two occupied residues.} Let the residues be $c,d$. Both multiplicities are at most $p+3<2p$, so $v^2=0$ for $v=y-z$. In addition $(z-c)(z-d)=0$. Hence
\begin{equation}\label{eq:two-only}
 z^2=(c+d)z-cd,\qquad
 yz=\tfrac12\bigl(y^2+(c+d)z-cd\bigr).
\end{equation}
The inverse of two exists because $p\ge5$.

\emph{Two thick residues and at least three occupied residues.} Write the thick multiplicities as $p+s_1,p+s_2$. If $s_1=s_2=1$, then, putting $u=x-z$ and $v=y-z$, both $uv=0$ and $v^2=0$. Thus
\begin{equation}\label{eq:two-thin}
 yz=xz+y^2-xy,\qquad z^2=2xz+y^2-2xy.
\end{equation}
Otherwise the budget forces the exact pattern $(p+2,p+1,1)$. Let $c$ be the residue of multiplicity $p+2$. Now $v^2=0$, $u^2v=0$, and $(z-c)uv=0$. Consequently
\begin{equation}\label{eq:three-relations}
 (y-z)^2=0,\qquad
 (x-c)(x-y)(y-z)=0,\qquad
 (y-c)(x-y)(y-z)=0.
\end{equation}
For the second identity, substitute $x-y=u-v$ to obtain $(z-c)uv+u^2v$, after terms containing $v^2$ vanish. The third identity is checked in the same way. The leading monomials in the three expanded identities are, respectively, $z^2$, $xyz$, and $y^2z$, with nonzero coefficients $1,1,1$. All other terms have smaller index. This supplies the stated reductions and exhausts the cases.
\end{proof}

\begin{corollary}[The upper bound at the new jump]\label{cor:upperjump}
For every prime $p\ge5$, $N_p^\circ(2p+4)\le p^2+p$.
\end{corollary}
\begin{proof}
Pad a smaller support with zero coefficients to obtain size $2p+4$. On a rational closed disk, a minimizing index above $p^2+p$ is above $t-1$, so Lemmas~\ref{lem:eligible} and~\ref{lem:cubic} exclude it. Any finite collection of roots in the open disk belongs to one rational closed disk. This proves the bound for the entire open disk, including multiplicities.
\end{proof}

\section{The two preceding endpoints}
\begin{lemma}[The endpoint cutoff]\label{lem:endpoint}
Let $t$ be $2p+2$ or $2p+3$, put $h=\lfloor(t-p)/2\rfloor$, and suppose $\sum_c\nu_c=t$. Every $T^k$ with $S_p(k)\le p-1$ and $k>p^2+h-1$ is reducible to earlier powers with no larger digit sum.
\end{lemma}
\begin{proof}
With no thick residue or one occupied residue the proof of Lemma~\ref{lem:cubic} applies immediately. Suppose there is one thick residue of multiplicity $p+s\le2p$, and let $m=t-p-s$ be the sum of the other multiplicities. Relations \eqref{eq:one-moderate} leave only monomials without $z$ or of the form $zx^a$. Translate the thick residue to zero. If $a\ge s$, then
\[
 zx^a=yx^a,
\]
because on the thick factor both sides vanish and on the others $y=z$. If $a\ge m$, division by the degree-$m$ defining polynomial of $R_{\rm other}$ expresses $zx^a$ in powers $zx^j$, $j<m$; multiplication by $z$ kills the thick factor. Thus every surviving $zx^a$ satisfies
\[
 a<\min(s,m)\le\lfloor(s+m)/2\rfloor=h.
\]

If the sole thick multiplicity is greater than $2p$, the others have total $m\le2$. In the translated coordinates $z^2=yz$, and $yz=zA(x)$ for a polynomial of degree less than $m$, because $y$ has such a representative in $R_{\rm other}$. Similarly $zx^a$ reduces once $a\ge m$. Thus a surviving exponent has $a\le1<h$.

It remains to treat two thick residues, with multiplicities $p+s_1,p+s_2$. Put $s=s_1+s_2\le3$. If there are exactly two occupied residues, \eqref{eq:two-only} eliminates $yz,z^2$. If there are three, the budget forces $s_1=s_2=1$ and one simple residue, so \eqref{eq:two-thin} does the same. For $v=y-z$, the polynomial
\[
 C(x)=(x-c_1)^{s_1}(x-c_2)^{s_2}
\]
annihilates $v$. If $a\ge s$, divide $x^a$ by $C$ and write its remainder as $A_a(x)$, of degree less than $s$. Then
\[
 zx^a=yx^a-A_a(x)y+A_a(x)z.
\]
This strictly reduces a monomial $zx^a$ without raising degree. Hence $a\le s-1\le2\le h-1$ for $p\ge5$. All reductions have the required digit-weight property, proving the lemma.
\end{proof}
The same minimizing-index argument now gives the upper bounds in \eqref{eq:first} and~\eqref{eq:second}.

\section{Attaining the new value}
Use the support \eqref{eq:support}, whose residue algebra is
\[
 R=\Fp[T]/T^{p+2}\ \times\ \Fp[U]/U^{p+1}\ \times\ \Fp,
\]
where $T$ is represented as $T$, $1+U$, and $2$ in the three factors. For an element $(f_0(T),f_1(U),f_2)$ define
\begin{align}
 \lambda(f_0,f_1,f_2)
 ={}&f_2-2f_1(0)+f_0(0)
 -2\bigl([U]f_1-[T]f_0\bigr)+2[T^2]f_0\notag\\
 &-\bigl([U^p]f_1-[T^p]f_0\bigr)+[T^{p+1}]f_0.
 \label{eq:lambda}
\end{align}
The coefficient notation always refers to the unique truncated representatives.

\begin{lemma}[The new independent moment]\label{lem:moment}
For $K=p^2+p$ one has $\lambda(T^K)=1$, while
\[
 \lambda(T^j)=0
 \quad\text{if }S_p(j)<2\text{ or if }S_p(j)=2\text{ and }j<K.
\]
\end{lemma}
\begin{proof}
A direct substitution into \eqref{eq:lambda} gives
\[
\begin{array}{c|rrrrrrrrr}
 f&1&x&y&z&x^2&xy&y^2&xz&yz\\\hline
 \lambda(f)&0&0&0&0&0&0&0&0&1.
\end{array}
\]
For instance $yz$ has local representatives $0$, $1+U^p$, and $4$, so its value is $4-2-1=1$. Every positive integer of digit sum one is $p^a$; for $a\ge2$ its class is $z$. The positive integers of digit sum two below $K$ are exactly $2,p+1,2p,p^2+1$, corresponding to the four displayed quadratic monomials before $yz$. The assertion for $j=0$ is included.
\end{proof}

Lift $\lambda$ to an integral linear functional $L$ on
\[
 \Z[T]\big/\prod_{e\in E}(T-e)
\]
by choosing integer lifts of its values on the power basis, with $L(1)=0$ exactly. All moments $M_j=L(T^j)$ are integral, and $M_K$ is a $p$-adic unit. Since the exponents in $E$ are distinct, Lagrange interpolation supplies unique rational coefficients $c_e$ satisfying $M_j=\sum_e c_ee^j$ for all $j$. Explicitly,
\begin{equation}\label{eq:lagrange}
 c_e=L\left(\prod_{f\in E\setminus\{e\}}\frac{T-f}{e-f}\right).
\end{equation}
Thus the functional is realized by a rational polynomial supported on $E$.

Choose $\delta$ as in \eqref{eq:depth}. The weight at $K$ is
\[
 w(K)=\frac2{p-1}+\delta K<\frac3{p-1}.
\]
An index of digit sum at least three has strictly larger weight. An index of digit sum two before $K$, or of digit sum one anywhere, has moment divisible by $p$ by Lemma~\ref{lem:moment}; its weight is at least $1+1/(p-1)>w(K)$. The constant moment is zero exactly, and later indices of digit sum two have larger weight by their index. Consequently $K$ is the unique minimizing index, and there are exactly $K$ roots in $\mathcal D(r)$ counted with multiplicities. The polynomial's degree is at most $K$, so it has degree exactly $K$ and all its roots lie in that disk.

To obtain integer coefficients, clear denominators. The strict coefficient gap on the chosen rational disk persists after sufficiently small $p$-adic perturbations of the allowed coefficients; this follows either from the convergent series or directly from the finite Taylor expansion at $X=1$. Fill any zero allowed coefficients by such rational perturbations. Finally vary the constant coefficient by a sufficiently small rational amount avoiding the finitely many critical values $-P(\alpha)$ with $P'(\alpha)=0$. This preserves the root count and makes every root simple. Avoiding one more value keeps the constant coefficient nonzero. Clearing denominators again gives exactly $2p+4$ nonzero integer coefficients. This proves the lower bound and attainment in \eqref{eq:third}.

\section{Attaining the preceding endpoints}
For completeness we include the earlier two-residue construction in the slight range extension needed here. Let $1\le h\le p-1$ and put
\[
 E_h=\{pj:0\le j<p+h\}\ \cup\ \{1+pj:0\le j<h\},\qquad K_h=p^2+h-1.
\]
The residue algebra is $\Fp[T]/T^{p+h}\times\Fp[U]/U^h$, with $T$ represented by $1+U$ in the second factor. Define
\[
 \lambda_h(f_0,f_1)
 =[U^{h-1}]f_1-[T^{h-1}]f_0-h[T^h]f_0-[T^{p+h-1}]f_0.
\]
Then
\[
 \lambda_h(T^j)=\binom{j}{h-1}-\mathbf1_{j=h-1}-h\mathbf1_{j=h}-\mathbf1_{j=p+h-1}
 \quad\text{in }\Fp.
\]
Since $h-1<p$, the binomial coefficient reduces to $\binom{a}{h-1}$ with $a=j\bmod p$. If $S_p(j)\le h$ and this is nonzero, the only possibilities are $j=h-1$, $j=h$, or $j=p^a+h-1$ for $a\ge1$. The three correction terms cancel exactly the possibilities of smaller digit sum, or equal digit sum and index below $K_h$. Therefore
\[
 \lambda_h(T^{K_h})=1,\qquad
 \lambda_h(T^j)=0\quad(S_p(j)<h\text{ or }S_p(j)=h,\ j<K_h),
\]
and $\lambda_h(1)=0$. Lift and interpolate as in \eqref{eq:lagrange}. For $0<\delta K_h<1/(p-1)$, the unit moment at $K_h$ has weight less than $(h+1)/(p-1)\le p/(p-1)$, whereas every obstructing positive-index moment is divisible by $p$ and has weight at least $p/(p-1)$. The constant moment vanishes exactly. Thus there are at least $K_h$ roots on the chosen disk, and they may be made simple by the perturbation argument above.

For $t=2p+2$ take $h=(p+1)/2$, which uses $p+2h=2p+1\le t$ terms. For $t=2p+3$ take $h=(p+3)/2$, which uses exactly $t$ terms and satisfies $h\le p-1$ because $p\ge5$. These lower bounds match Lemma~\ref{lem:endpoint}, completing Theorem~\ref{thm:main}.

\section{Computational evidence and scope}
The finite computations are checks of the proofs, not assumptions of the uniform theorem. The script \texttt{new\_phase\_checks.py} constructs the functional witnesses using exact rational arithmetic, clears denominators, perturbs them if needed, and checks the last minimum of the integer Taylor coefficient valuations at $X=1$. It also verifies squarefreeness by an exact polynomial gcd. The companion \texttt{modp\_scan.py} enumerates residue multiplicity profiles and their digit-weight filtrations.

\begin{center}
\begin{tabular}{rrrrr}\toprule
$p$&nonzero terms&degree&roots in the disk&closed depth used\\\midrule
5&14&30&30&$61/240$\\
7&18&56&56&$113/672$\\
11&26&132&132&$265/2640$\\
13&30&182&182&$365/4368$\\
17&38&306&306&$613/9792$\\\bottomrule
\end{tabular}
\end{center}
These depths equal $1/(p-1)+1/(2(p-1)(p^2+p))$. All displayed witnesses are squarefree. The complete coefficients are in \texttt{new\_phase\_certificates.json}. Three exhaustive residue-profile scans for $p=5$, after placing a residue of maximal multiplicity at zero, check $419$, $541$, and $691$ profiles at $t=12,13,14$ and yield $27,28,30$. A further $21$ exact local-algebra checks cover the distinct reduction patterns at $p=5,7,11$. The modular scan is independently supported by the integer Taylor-coefficient checks of the attaining polynomials; neither is used to infer the infinite-prime theorem.

\paragraph{Relation to the source program.}
The input is the ProveIt GRH article's use of integral recurrences, residue algebras, valuation comparisons and sparse polynomial obstructions, specifically its singular-moduli section. No hypothesis about zeros of an $L$-function occurs in this note. The latest report \cite{transport} is expressly credited for the two-residue construction and the previously proved range; the new contribution is the upper cutoff at $2p+4$, its three-residue attaining functional, and the two adjacent endpoint values. The present note does not settle Riffaut's full sparse singular-modulus conjecture or a classical RH-type problem.

\paragraph{Literature and verification status.}
Classical work on $p$-adic exponential polynomials gives bounds in terms of their order \cite{vpr}, and more recent sparse-polynomial research treats algorithmic root counts and root separation over specified local fields \cite{rz}. The exact threshold proved here concerns geometric roots in a specified deep disk and the indicated growing term budget. Targeted searches found no matching statement, but they do not establish worldwide priority. The proof is human-readable and independently checkable; no compiled Lean formalization is claimed.

\clearpage
\begin{thebibliography}{9}
\raggedright
\bibitem{article} ProveIt project, \emph{Beyond the Riemann hypothesis for Dirichlet $L$-functions: consequences, transfers and limits}, section \texttt{smd.tex}, snapshot \texttt{ce35f320bc88c02093e479f5fe5003ff86b2795f}. \url{https://github.com/VladimirReshetnikov/ProveIt/tree/ce35f320bc88c02093e479f5fe5003ff86b2795f/Analysis/GeneralizedRiemannHypothesis/Article}.
\bibitem{transport} ProveIt incoming report, \emph{Exact residue transport for sparse $p$-adic roots: arbitrary supports, exponential polynomials, and the range from $p+2$ to $2p+1$ terms}, 27 November 2026, delivered as \texttt{padic\_sparse\_research.zip}; Theorems 1--3 and the two-residue functional in the proof of the sharp range. Inspected in the same repository snapshot.
\bibitem{vpr} A.~J. van der Poorten and R.~S. Rumely, \emph{Zeros of $p$-adic exponential polynomials II}, Journal of the London Mathematical Society (2) \textbf{36} (1987), 1--15. \url{https://doi.org/10.1112/jlms/s2-36.1.1}.
\bibitem{rz} J.~M. Rojas and Y. Zhu, \emph{Root repulsion and faster solving for very sparse polynomials over $p$-adic fields}, Journal of Number Theory \textbf{241} (2022), 655--699. \url{https://doi.org/10.1016/j.jnt.2022.01.013}; preprint \url{https://arxiv.org/abs/2107.09173}.
\end{thebibliography}
\end{document}
